\documentclass[10pt,aspectratio=169]{beamer}

\usetheme{Madrid}
\usecolortheme{dolphin}
\usefonttheme{professionalfonts}

\usepackage[T1]{fontenc}
\usepackage[utf8]{inputenc}
\usepackage{amssymb}
\usepackage{booktabs}
\usepackage{xcolor}
\usepackage{tikz-cd}

\definecolor{deepblue}{RGB}{0,70,130}
\setbeamercolor{title}{fg=deepblue}
\setbeamercolor{frametitle}{fg=deepblue}
\setbeamertemplate{navigation symbols}{}
\setbeamertemplate{itemize items}[circle]

\newcommand{\code}[1]{\texttt{\small\detokenize{#1}}}

% Map macros, matching the manuscript (sections/chyll_v2.tex).
\newcommand{\Enc}{E_\theta}
\newcommand{\Dec}{D_\theta}
\newcommand{\Dyn}{F_\theta}
\newcommand{\Lat}{Z}
\newcommand{\Xr}{X'}

\title[CHyLL v2]{CHyLL v2: Continuous Latent Embeddings of Hybrid Systems}
\subtitle{What the model learns, and what each loss term enforces}
\author{Kaito Iwasaki}
\date{}

\begin{document}

\begin{frame}
  \titlepage
\end{frame}

\begin{frame}{Hybrid Dynamical Systems}
  A \emph{hybrid} dynamical system mixes smooth motion with sudden jumps.

  \vspace{0.8em}
  \begin{itemize}
    \item A walking robot: the swing leg moves smoothly, then strikes the
      ground and the state jumps.
    \item A bouncing ball: free flight under gravity, then an instantaneous
      change of velocity at each impact.
  \end{itemize}

  \vspace{0.8em}
  The jump is a genuine discontinuity. A smooth model of the dynamics cannot
  be fitted directly across it.
\end{frame}

\begin{frame}{Three Example Systems}
  \begin{center}
  \begin{tabular}{lcl}
    \toprule
    System & State dim. & Behavior \\
    \midrule
    Rimless wheel & 2 & stable rolling limit cycle \\
    Bouncing ball & 2 & bounces, losing energy at each impact \\
    Compass-gait biped & 4 & period-doubling route to chaos \\
    \bottomrule
  \end{tabular}
  \end{center}

  \vspace{1em}
  They increase in difficulty. The compass-gait biped walks with period 1,
  then 2, 4, 8, and finally chaotically as the ground slope steepens.
\end{frame}

\begin{frame}{The Idea}
  Replace the discontinuous system with a continuous one that is easy to model.

  \vspace{0.8em}
  \begin{itemize}
    \item An \textbf{encoder} maps each state into a latent space.
    \item A single \textbf{smooth vector field} advances the latent state in
      time.
    \item A \textbf{decoder} maps the latent state back.
  \end{itemize}

  \vspace{0.8em}
  The jump is absorbed into the smooth latent flow. This follows Teng, Liu,
  and Sreenath (ICML 2026).
\end{frame}

\begin{frame}{The Mapping-Cylinder State Space}
  To make the jump continuous, the state is augmented with one extra
  coordinate \(s \in [0,1]\):
  \[
    y = (x, s) \in X'.
  \]

  \begin{itemize}
    \item \(s = 0\): ordinary motion of the original system.
    \item \(s\) sweeping \(0 \to 1\): the jump, unrolled as a continuous path.
  \end{itemize}

  \vspace{0.6em}
  On \(X'\) the trajectory is continuous, so a smooth latent model can be fitted
  to it.
\end{frame}

\begin{frame}{The Learned Maps}
  CHyLL v2 learns three maps. Write \(Z = \mathbb{R}^d\) for the latent space.

  \vspace{0.6em}
  \begin{align*}
    E_\theta &: X' \to Z && \text{encoder},\\
    V_\theta &: Z \to Z && \text{latent vector field},\\
    D_\theta &: Z \to X' && \text{decoder}.
  \end{align*}

  \vspace{0.4em}
  The time-\(\tau\) flow of \(V_\theta\) is written \(F_\theta\). The latent
  dimension is \(d = 7\) for the rimless wheel and bouncing ball, and
  \(d = 11\) for the compass-gait biped.
\end{frame}

\begin{frame}{Training: The Forward Pass}
  For a sampled trajectory slice \(y_0, \ldots, y_{T-1}\):

  \vspace{0.5em}
  \begin{align*}
    z_i &= E_\theta(y_i),\\
    \widehat z_i &= F_\theta^{\,i}(z_0),\\
    \widehat y_i &= D_\theta(\widehat z_i).
  \end{align*}

  \vspace{0.4em}
  Each state is encoded, the latent flow is integrated forward, and the
  integrated latent trajectory is decoded. The loss compares the encoded
  \(z_i\), the integrated \(\widehat z_i\), and the decoded \(\widehat y_i\)
  against the data.
\end{frame}

\begin{frame}[fragile]{The Commutative Diagram}
  Here \(f_\tau\) is one time-\(\tau\) step of the system, \(G\) the impact
  set, \(r\) the reset map, and \(\iota_0, \iota_1\) the inclusions at
  \(s = 0\) and \(s = 1\). The learned maps should make this diagram commute:

  \vspace{0.3em}
  \begin{center}
  \small
  \begin{tikzcd}[column sep=3.0em, row sep=2.4em]
    G \arrow[r, "\iota_1"] \arrow[d, "\iota_0 \circ r"']
    & \Xr \arrow[r, "f_\tau"] \arrow[d, "\Enc"']
    & \Xr \arrow[d, "\Enc"] \\
    \Xr \arrow[r, "\Enc"'] \arrow[rd, "\mathrm{id}_{\Xr}"']
    & \Lat \arrow[r, "\Dyn"'] \arrow[d, "\Dec"']
    & \Lat \\
    & \Xr
    \arrow[phantom, from=2-1, to=3-2]
    &
  \end{tikzcd}
  \end{center}

  \vspace{0.2em}
  \normalsize
  Each cell is enforced by one loss term:
  \begin{itemize}
    \item Left square, \(L_g\): \(E_\theta\) sends both sides of a reset to one latent point.
    \item Top-right square, \(L_z\): the latent flow matches the dynamics.
    \item Lower triangle, \(L_x\): the decoder inverts the encoder.
  \end{itemize}
\end{frame}

\begin{frame}{Two Conditions Beyond the Diagram}
  Commuting is not the whole story. Two more terms act at the seam where
  \(s\) wraps from \(1\) back to \(0\):

  \vspace{0.6em}
  \begin{itemize}
    \item \(L_v\): the latent velocity has no kink crossing a seam.
    \item \(L_c\): the encoder keeps all \(d\) latent coordinates active and
      does not collapse them.
  \end{itemize}

  \vspace{0.6em}
  These seams are detected from the sampled coordinate \(s\) in the data, not
  from symbolic guard equations.
\end{frame}

\begin{frame}{What Each Loss Term Enforces}
  \small
  \begin{center}
  \begin{tabular}{lll}
    \toprule
    Term & Enforces & Code helper \\
    \midrule
    \(L_x\) & decoder recovers the state & \code{reconstruction_loss} \\
    \(L_z\) & encode/flow square commutes & \code{latent_consistency_loss} \\
    \(L_g\) & reset pairs share a latent point & \code{gluing_loss} \\
    \(L_v\) & latent velocity continuous at seams & \code{seam_velocity_loss} \\
    \(L_c\) & latent coordinates do not collapse & \code{anti_collapse_loss} \\
    \bottomrule
  \end{tabular}
  \end{center}

  \vspace{0.8em}
  \normalsize
  \(L_g\), \(L_z\), and \(L_x\) are the three commuting cells of the diagram.
  \(L_v\) and \(L_c\) are the two conditions beyond it.
\end{frame}

\begin{frame}{The Composite Objective}
  \[
    L = w_x L_x + w_z L_z + w_g L_g + w_v L_v + w_c L_c.
  \]

  \vspace{0.8em}
  \begin{center}
  \begin{tabular}{lc}
    \toprule
    Setting & \((w_x, w_z, w_g, w_v, w_c)\) \\
    \midrule
    Rimless wheel, bouncing ball & \((1,\ 1,\ 3,\ 1,\ 1)\) \\
    Compass-gait biped & \((1,\ 1,\ 3,\ 0,\ 1)\) \\
    \bottomrule
  \end{tabular}
  \end{center}

  \vspace{0.8em}
  The exact values for any run are stored in its \code{config.json}.
\end{frame}

\begin{frame}{Results}
  \begin{itemize}
    \item \textbf{Rimless wheel.} The model reproduces the rolling limit cycle,
      and the learned embedding recovers its expected single topological cycle.
    \item \textbf{Compass-gait biped.} The cascade analysis recovers the
      period-doubling sequence exactly: the learned latent states form 2, 4,
      and 8 clusters at the period-2, -4, and -8 gaits, and disperse under
      chaos.
    \item \textbf{Bouncing ball.} The orbit spirals slowly toward rest rather
      than closing into a loop, so its topological signature is only partially
      recovered.
  \end{itemize}
\end{frame}

\begin{frame}{Where the Code Lives}
  \small
  \begin{center}
  \begin{tabular}{ll}
    \toprule
    Object & File \\
    \midrule
    \(E_\theta, V_\theta, D_\theta\) & \code{networks.py} \\
    latent ODE rollout & \code{ode.py} \\
    loss terms \(L_x, L_z, L_g, L_v, L_c\) & \code{losses.py} \\
    training loop & \code{train.py} \\
    system simulators and reset maps & \code{systems/} \\
    \bottomrule
  \end{tabular}
  \end{center}

  \vspace{0.7em}
  \normalsize
  These files are the CHyLL v2 core library, under \code{chyll_v2/chyll_v2/}.
  Topology postprocessing lives under \code{chyll_v2/cycling_signature/}.
\end{frame}

\begin{frame}{Summary}
  \begin{itemize}
    \item A hybrid system is made continuous on the mapping cylinder \(X'\).
    \item CHyLL v2 learns an encoder, a latent vector field, and a decoder.
    \item \(L_g\), \(L_z\), and \(L_x\) are the commuting cells of the diagram.
    \item \(L_v\) and \(L_c\) add the seam-regularity conditions.
    \item The learned embeddings recover the expected topology for the rimless
      wheel and the compass-gait cascade.
  \end{itemize}
\end{frame}

\end{document}
